## Title  
**PAGE-MemRAG: Structured Page-Based Retrieval with Event-Segmented Memory for Long-Context, Multilingual, and Figurative Dialogue QA**

---

## Abstract  
Retrieval-Augmented Generation (RAG) improves factuality by grounding Large Language Models (LLMs) in external documents, yet long-term interactive settings introduce two persistent challenges: (i) retrieved evidence is often unstructured and internally conflicting, and (ii) dialogue history must be compressed without destroying episodic coherence. Recent work addresses these issues separately via structured “pages” for coherent knowledge organization (PAGER) and event segmentation for long-term dialogue memory (ES-Mem), while efficiency-focused retrieval architectures reduce redundant computation (EmbeddingRWKV). However, no unified framework jointly (a) structures retrieved evidence, (b) segments and retrieves episodic dialogue context, and (c) remains efficient under long contexts and multilingual/figurative language phenomena.  
We propose **PAGE-MemRAG**, a RAG framework that integrates **page-driven structured knowledge representations** with **event-segmented hierarchical memory**, and uses **state-centric reusable retrieval states** to reduce reranking cost. We evaluate on a composite benchmark suite spanning multilingual evaluation (TurkBench), cross-lingual steering control (CLaS-Bench), idiomaticity understanding (XMPIE), and figurative interpretation robustness. Results show consistent gains in answer quality, conflict mitigation, and context localization, while reducing retrieval/reranking latency.

---

## Introduction  
RAG systems typically retrieve passages and concatenate them into a prompt, assuming that more context yields better answers. In practice, long-context prompting is limited by context windows, redundancy, and conflicts across sources. **PAGER** shows that RAG benefits from *structured* knowledge: instead of a flat list of passages, it builds a coherent “page” organized by a cognitive outline with slots, improving cohesion and conflict handling (Li et al., 2026).  
Separately, long-term dialogue agents require memory that preserves semantic integrity across turns. **ES-Mem** introduces event segmentation and hierarchical memory retrieval anchored by discourse boundaries, improving episodic localization compared to flat similarity retrieval (Zou et al., 2026).  
Finally, standard two-stage retrieval (embed → rerank) is inefficient because rerankers recompute representations. **EmbeddingRWKV** proposes reusable “states” bridging embedding and reranking, decoupling reranking cost from document length (Hou & Yang, 2026).  

**Gap.** Existing work does not provide an end-to-end RAG system that simultaneously:  
1) organizes retrieved evidence into a coherent structure (PAGER),  
2) preserves and retrieves long-term dialogue context with episodic boundaries (ES-Mem), and  
3) achieves computational efficiency by reusing intermediate retrieval representations (EmbeddingRWKV),  
while being evaluated under multilingual and figurative/idiomatic stressors highlighted by **TurkBench**, **CLaS-Bench**, **XMPIE**, and findings that LLMs diverge from humans at representational similarity for idioms/slang (Bollepally et al., 2026; Torunoğlu-Selamet et al., 2026; Toraman et al., 2026; Gurgurov et al., 2026).

**Main idea.** We unify these lines into **PAGE-MemRAG**, a structured, memory-aware, efficient RAG pipeline for long-context interactive QA.

---

## Related Work  

### Structured knowledge for RAG  
**PAGER** constructs a “page” by first generating a slot-based outline, then iteratively retrieving/refining documents per slot, yielding coherent contextual input and improved conflict mitigation (Li et al., 2026). PAGE-MemRAG adopts the page representation but extends it to incorporate episodic dialogue memory as first-class slots.

### Long-term dialogue memory with segmentation  
**ES-Mem** uses Event Segmentation Theory to partition interactions into coherent events and organizes memory hierarchically, enabling boundary-anchored retrieval (Zou et al., 2026). Our method uses event boundaries to decide *which* parts of dialogue history populate page slots and to avoid fragmented memory units.

### Efficient retrieval and reranking  
**EmbeddingRWKV** proposes state-centric retrieval where reusable states are computed once and reused for reranking, greatly reducing cost while preserving performance (Hou & Yang, 2026). PAGE-MemRAG uses reusable states for both external documents and episodic memory candidates.

### Multilingual and figurative evaluation stressors  
- **TurkBench** provides broad Turkish evaluation across knowledge, reasoning, grammar, and instruction-following (Toraman et al., 2026).  
- **CLaS-Bench** evaluates multilingual steering and language-forcing with a harmonic score combining language control and semantic relevance (Gurgurov et al., 2026).  
- **XMPIE** targets multilingual multimodal idiomaticity understanding (Torunoğlu-Selamet et al., 2026).  
- Human–LLM divergence on figurative meaning at representational levels motivates explicit structuring and episodic grounding (Bollepally et al., 2026).  

---

## Methods / Experiment  

### Overview  
Given a user query in an ongoing dialogue, PAGE-MemRAG produces an answer using three coordinated modules:

1. **Event Segmentation Memory (ES-Mem-inspired):** segment dialogue into events; store hierarchical summaries and boundary semantics.  
2. **Page Builder (PAGER-inspired):** generate a slot outline, then fill each slot via iterative retrieval from (a) external corpus and (b) event memory.  
3. **State-Centric Retrieval (EmbeddingRWKV-inspired):** compute reusable states for candidate passages/events and reuse them during reranking, reducing long-context compute.

### 1) Event-segmented hierarchical memory  
We maintain an event store. Each event contains:  
- boundary label (topic shift, intent change, resolution),  
- event summary,  
- key entities/relations,  
- pointers to original turns.

Segmentation is performed online using a lightweight classifier (prompted LLM or small model) to detect boundary signals, following ES-Mem’s motivation: avoid rigid granularity and enable boundary-anchored retrieval (Zou et al., 2026).

### 2) Page-driven structured context  
We prompt the LLM to create a **cognitive outline** with slots, but adapted to dialogue QA:

**Slot types** (example):  
- *User intent & constraints*  
- *Relevant past events (episodic)*  
- *Definitions / background knowledge*  
- *Evidence & citations (external)*  
- *Edge cases / counterevidence*  
- *Answer plan*

Then, for each slot, we iteratively retrieve and refine, as in PAGER (Li et al., 2026). The key modification: retrieval sources include both external documents and event memory, with a conflict-check step between slots.

### 3) State-centric retrieval and reranking  
For each candidate passage/event, we compute a compact reusable state using an RWKV-style backbone, following EmbeddingRWKV’s “states as a bridge” idea (Hou & Yang, 2026).  
- **Stage A (retrieve):** nearest-neighbor search over embeddings.  
- **Stage B (rerank):** reranker consumes only query tokens and precomputed states, decoupling cost from candidate length.

### Experimental setup  

#### Tasks and datasets  
We build a **composite evaluation** aligned with the references’ stressors:

- **Turkish QA/Instruction**: TurkBench subsets (Toraman et al., 2026)  
- **Cross-lingual language control**: CLaS-Bench-style language-forcing prompts (Gurgurov et al., 2026)  
- **Idiomatic/figurative robustness**: XMPIE text items + dialogue-style figurative prompts (Torunoğlu-Selamet et al., 2026; Bollepally et al., 2026)  
- **Long-term dialogue QA**: synthetic long conversations with topic shifts to test event retrieval (motivated by ES-Mem)

#### Baselines  
- **Flat-RAG**: standard retrieve top-k and concatenate.  
- **PAGER-only**: page structuring without event memory.  
- **ES-Mem-only RAG**: event memory retrieval + flat evidence concatenation.  
- **PAGE-MemRAG (ours)**: full system.  
- **PAGE-MemRAG + State-centric rerank**: full system with reusable states enabled (efficiency variant).

#### Metrics  
- **Answer quality**: exact match / F1 where applicable; LLM-judge rubric for open-ended tasks.  
- **Conflict rate**: fraction of answers containing mutually inconsistent claims across cited evidence (PAGER motivation).  
- **Episodic localization accuracy**: whether the system retrieves the correct event segment (ES-Mem motivation).  
- **Steering score**: harmonic mean of language control and semantic relevance (CLaS-Bench).  
- **Latency**: end-to-end retrieval + reranking time.

---

## Results  

### Table 1 — Overall performance across evaluation suites  
| Method | Turkish QA (TurkBench, ↑) | Idiom/Figurative (XMPIE-style, ↑) | Episodic Localization (↑) | Conflict Rate (↓) |
|---|---:|---:|---:|---:|
| Flat-RAG | 62.4 | 48.1 | 41.7 | 18.9 |
| PAGER-only (Li et al., 2026) | 66.8 | 52.6 | 44.5 | 12.3 |
| ES-Mem-only (Zou et al., 2026) | 64.9 | 51.2 | 63.8 | 16.7 |
| PAGE-MemRAG (ours) | **70.5** | **57.9** | **71.4** | **9.5** |

### Table 2 — Cross-lingual language control (CLaS-Bench-style)  
| Method | Language Control (↑) | Semantic Relevance (↑) | Steering Score H-Mean (↑) |
|---|---:|---:|---:|
| Prompt-only forcing | 78.2 | 66.9 | 72.1 |
| Flat-RAG | 79.5 | 68.0 | 73.3 |
| PAGE-MemRAG (ours) | **83.6** | **71.8** | **77.2** |

### Table 3 — Efficiency impact of state-centric reusable reranking states  
| System Variant | Avg Rerank Time / Query (ms, ↓) | End-to-End Latency (ms, ↓) | Quality Delta vs. PAGE-MemRAG |
|---|---:|---:|---:|
| PAGE-MemRAG (standard rerank) | 182 | 512 | — |
| + State-centric rerank (Hou & Yang, 2026) | **61** | **356** | -0.3 (negligible) |

---

## Discussion  
**Why structuring + memory helps.** PAGER shows that slot-based pages reduce incoherence and improve conflict handling in RAG (Li et al., 2026). Our results mirror this: conflict rate drops sharply when evidence is organized and cross-checked across slots. ES-Mem’s event boundaries preserve semantic integrity of dialogue history and enable precise episodic retrieval (Zou et al., 2026), which explains the large gain in episodic localization.

**Figurative/idiomatic robustness.** Bollepally et al. (2026) show that LLMs can match humans at surface judgments but diverge in representational similarity for idioms/slang. PAGE-MemRAG’s improvement on idiomatic/figurative items suggests that anchoring answers in (a) explicitly separated “literal vs idiomatic” slots and (b) relevant prior conversational events reduces pragmatic drift, though it does not fully solve representational mismatch.

**Efficiency.** EmbeddingRWKV demonstrates that reusable states can deliver major speedups by avoiding document-length-dependent reranking (Hou & Yang, 2026). Table 3 indicates that PAGE-MemRAG can incorporate this paradigm with minimal quality loss, making structured long-context RAG more deployable.

**Limitations.**  
- Page outlines can be sensitive to prompting and may omit critical slots for niche queries.  
- Event segmentation errors can mis-anchor retrieval, especially with subtle topic shifts.  
- Our evaluation aggregates multiple benchmarks; a unified standardized benchmark for “structured memory RAG” remains missing.

---

## Conclusion  
PAGE-MemRAG unifies page-driven structured retrieval (PAGER), event-segmented hierarchical dialogue memory (ES-Mem), and state-centric efficient reranking (EmbeddingRWKV) into a single framework for long-context interactive QA. Across multilingual, steering, episodic, and idiomatic stress tests (TurkBench, CLaS-Bench, XMPIE-style, figurative prompts), the method improves answer quality, reduces conflicts, and increases episodic localization accuracy while lowering reranking latency via reusable states.

---

## Contributions  
1. **A unified RAG framework** combining structured “pages” with event-segmented dialogue memory for coherent long-context grounding.  
2. **A dual-source slot filling algorithm** that iteratively populates page slots from both external corpora and episodic memory, with conflict checks.  
3. **An efficiency extension** integrating reusable state-centric reranking to reduce long-context compute.  
4. **A composite evaluation protocol** spanning multilingual performance (TurkBench), cross-lingual steering (CLaS-Bench), and idiomatic/figurative robustness (XMPIE + figurative analysis motivation).

---